\chapter*{Abstract}

In many settings, a query's result must stay current as new data arrives. Incremental view maintenance (IVM) provides this by maintaining a structure that encodes the result and updating it from each change, rather than recomputing it from scratch. Abo~Khamis et al.\ give an approach, which we call IVM$^+$, that maintains full conjunctive queries in the insert-only setting by organising the computation around a tree decomposition of the query, with a cost governed by the query's fractional hypertree width. Although the approach is theoretical, in the insert-only setting it can be expressed entirely as a set of SQL views, and so realised on top of an existing incremental engine without modifying it.

This thesis investigates whether that realisation is beneficial in practice. We build a query decomposer that rewrites a join query into the bag views of a tree decomposition, in both IVM$^+$ and an adjusted variant that drops its semijoin filters, together with a semiring extension for computing aggregates such as \texttt{COUNT}. We evaluate the rewrites on Feldera, a general-purpose incremental engine, against the baseline of maintaining the query directly, without decomposition, across graph, TPC-H and IMDB workloads.

The rewrite pays off on the queries the baseline handles worst, cyclic queries over dense data: there it maintains one query $42.96$ times faster than the baseline, at the same memory, and keeps tractable several queries the baseline cannot maintain within three hours. These gains hold even though the engine does not implement the worst-case-optimal joins the width bound assumes. Where the baseline is already cheap, on acyclic key-based joins and sparse data, the rewrite gives no speedup and adds a memory overhead of up to $2.6$ times. Within the decomposition, the final reconstruction dominates the cost and is avoidable when only the result's size is needed. Realised as plain SQL views on an unmodified engine, the IVM$^+$ decomposition is thus a worthwhile rewrite for cyclic queries over dense data, and query-shape-dependent in general.
